%% mwe_creditos_titulo_invisivel.tex
%% Pergunta (P11 do plan-est_2026, segunda revisão de 24/09/2026):
%%   o modelo de créditos do Relatório de Gestão pode ir para a linha
%%   relatorio (PDF/UA-2) com a página de créditos SEM título impresso
%%   ("Créditos institucionais" etc.), mas com o título na árvore de
%%   tags, para o leitor de tela navegar?
%%
%% Variantes (por \def\variante na linha de comando):
%%   A  título invisível com texto real: \@startsection de pele zero,
%%      texto em Tr 3 (pdfrender), \smash\rlap, \topskip emulado.
%%   B  título sem conteúdo: \tagstructbegin{tag=H1,alt=...} vazio.
%%   C  sem título nenhum — a referência visual.
%% E, nas três, o que o modelo do RG pede além do título:
%%   - Plenário em duas colunas (o \parbox duplo do RG) pelo
%%     NomesDuasColunas do book;
%%   - \NomeNosCreditos e \CargoNosCreditos na interface do book 0.5.1,
%%     com o registro separado do nome (decisão de 17/09);
%%   - o fio sob o título de nível 2 (\CreditoInstitucionalUm) como
%%     artefato.
%%
%% RESULTADO (2026-09-24, TL2026, lualatex, veraPDF ua2):
%%   A  PASS. Árvore certa: Sect > H1 "Créditos institucionais" > Sect
%%      > H2 … ; o /T do H1 é o texto limpo; pdftotext extrai o título.
%%      A página sai IDÊNTICA à de C, pixel a pixel (pdftoppm -r 60 +
%%      diferença de imagem vazia), e o título seguinte cai em
%%      97,3956 pt nas duas. → É o padrão a portar.
%%   B  PASS no veraPDF, mas a estrutura quebra: o H1 fica solto, sem
%%      Sect, e o da 2ª página cai DENTRO do Sect de "Comissões
%%      permanentes" da 1ª. Um H1 sem conteúdo também depende de o
%%      leitor de tela ler /Alt em título, o que não é garantido.
%%      → Descartada.
%%   Nomes e cargos: interface do book 0.5.1 (book-crpsp_acessivel.sty
%%      :781-810) com a decisão de 17/09 aplicada — registro separado
%%      do nome, com o nome \NomeNosCreditos mantido (editor, 24/09):
%%        \NomeNosCreditos{nome}{registro}       forma nova
%%        \NomeNosCreditos{nome (CRP~06/nº)}     antiga, com aviso
%%        \CargoNosCreditos[registro]{nome}{cargo}
%%      Árvore: um P por pessoa, com o registro num Span próprio
%%      (14 Span no MWE: 10 nomes + 4 da Diretoria). A forma antiga sai
%%      como antes, nome e registro no mesmo MC.
%%   Duas colunas: NomesDuasColunas do book (multicols, \\ vira \par).
%%      Um P por pessoa, na ordem coluna 1 → coluna 2. O registro vai
%%      sempre para a linha de baixo, recuado (aprovado pelo editor em
%%      24/09). Não precisa do padrão de tagging do `paralelos'.
%%   Espaço entre pessoas: \crpspCreditosSep (0,3 \baselineskip) no
%%      CreditoInstitucional e no NomesDuasColunas. É exceção ao \parskip
%%      do WCAG numa lista de nomes — A CONFIRMAR COM O EDITOR.
%%      Do título ao primeiro nome: 5,67 pt nos dois ambientes.
%%   Porte (plan-est_2026, comum.tex §5, 24/09): acrescentou a nota
%%      opcional \NomeNosCreditos{nome}{registro}[licenciada] e a
%%      proibição de partir uma pessoa entre colunas/páginas (\clubpenalty,
%%      \widowpenalty e \interlinepenalty 10000 nos dois ambientes) —
%%      no pe26 o registro de uma pessoa foi parar na página seguinte.
%%   E  (terceira revisão, 24/09) cargos alinhados: \CargoNosCreditos
%%      com o nome numa \parbox de largura fixa (\crpspCreditosNome,
%%      0,56 \linewidth) e o cargo noutra, com o tagging do `paralelos'.
%%      PASS; por pessoa, Div > Div(P nome + Span registro) + Div(P
%%      cargo), na ordem certa, sem P vazio. O registro vai sempre na
%%      linha de baixo: na coluna estreita ele ora cabia, ora não.
%%   Fio: fora da árvore (nenhum nó novo), sem conteúdo não marcado.
%%   Armadilhas achadas no caminho estão em comentário junto de cada
%%   macro: \rule num parágrafo com o tagging suspenso desbalanceia os
%%   hooks; o \vspace{-\parskip} da crpsp-base puxa o fio para cima do
%%   título, o multicols e o primeiro nome para dentro dele; o tipo `g'
%%   do ltcmd saiu do kernel.
%%
%% Compilar (2 passadas por variante):
%%   export TEXINPUTS="$HOME/repos/crpsp-latex/crpsp-book/desenvolvimento/v2:"
%%   for v in A B C E; do for i in 1 2; do lualatex -jobname=mwe_creditos_$v \
%%     "\def\variante{$v}%% mwe_creditos_titulo_invisivel.tex
%% Pergunta (P11 do plan-est_2026, segunda revisão de 24/09/2026):
%%   o modelo de créditos do Relatório de Gestão pode ir para a linha
%%   relatorio (PDF/UA-2) com a página de créditos SEM título impresso
%%   ("Créditos institucionais" etc.), mas com o título na árvore de
%%   tags, para o leitor de tela navegar?
%%
%% Variantes (por \def\variante na linha de comando):
%%   A  título invisível com texto real: \@startsection de pele zero,
%%      texto em Tr 3 (pdfrender), \smash\rlap, \topskip emulado.
%%   B  título sem conteúdo: \tagstructbegin{tag=H1,alt=...} vazio.
%%   C  sem título nenhum — a referência visual.
%% E, nas três, o que o modelo do RG pede além do título:
%%   - Plenário em duas colunas (o \parbox duplo do RG) pelo
%%     NomesDuasColunas do book;
%%   - \NomeNosCreditos e \CargoNosCreditos na interface do book 0.5.1,
%%     com o registro separado do nome (decisão de 17/09);
%%   - o fio sob o título de nível 2 (\CreditoInstitucionalUm) como
%%     artefato.
%%
%% RESULTADO (2026-09-24, TL2026, lualatex, veraPDF ua2):
%%   A  PASS. Árvore certa: Sect > H1 "Créditos institucionais" > Sect
%%      > H2 … ; o /T do H1 é o texto limpo; pdftotext extrai o título.
%%      A página sai IDÊNTICA à de C, pixel a pixel (pdftoppm -r 60 +
%%      diferença de imagem vazia), e o título seguinte cai em
%%      97,3956 pt nas duas. → É o padrão a portar.
%%   B  PASS no veraPDF, mas a estrutura quebra: o H1 fica solto, sem
%%      Sect, e o da 2ª página cai DENTRO do Sect de "Comissões
%%      permanentes" da 1ª. Um H1 sem conteúdo também depende de o
%%      leitor de tela ler /Alt em título, o que não é garantido.
%%      → Descartada.
%%   Nomes e cargos: interface do book 0.5.1 (book-crpsp_acessivel.sty
%%      :781-810) com a decisão de 17/09 aplicada — registro separado
%%      do nome, com o nome \NomeNosCreditos mantido (editor, 24/09):
%%        \NomeNosCreditos{nome}{registro}       forma nova
%%        \NomeNosCreditos{nome (CRP~06/nº)}     antiga, com aviso
%%        \CargoNosCreditos[registro]{nome}{cargo}
%%      Árvore: um P por pessoa, com o registro num Span próprio
%%      (14 Span no MWE: 10 nomes + 4 da Diretoria). A forma antiga sai
%%      como antes, nome e registro no mesmo MC.
%%   Duas colunas: NomesDuasColunas do book (multicols, \\ vira \par).
%%      Um P por pessoa, na ordem coluna 1 → coluna 2. O registro vai
%%      sempre para a linha de baixo, recuado (aprovado pelo editor em
%%      24/09). Não precisa do padrão de tagging do `paralelos'.
%%   Espaço entre pessoas: \crpspCreditosSep (0,3 \baselineskip) no
%%      CreditoInstitucional e no NomesDuasColunas. É exceção ao \parskip
%%      do WCAG numa lista de nomes — A CONFIRMAR COM O EDITOR.
%%      Do título ao primeiro nome: 5,67 pt nos dois ambientes.
%%   Porte (plan-est_2026, comum.tex §5, 24/09): acrescentou a nota
%%      opcional \NomeNosCreditos{nome}{registro}[licenciada] e a
%%      proibição de partir uma pessoa entre colunas/páginas (\clubpenalty,
%%      \widowpenalty e \interlinepenalty 10000 nos dois ambientes) —
%%      no pe26 o registro de uma pessoa foi parar na página seguinte.
%%   E  (terceira revisão, 24/09) cargos alinhados: \CargoNosCreditos
%%      com o nome numa \parbox de largura fixa (\crpspCreditosNome,
%%      0,56 \linewidth) e o cargo noutra, com o tagging do `paralelos'.
%%      PASS; por pessoa, Div > Div(P nome + Span registro) + Div(P
%%      cargo), na ordem certa, sem P vazio. O registro vai sempre na
%%      linha de baixo: na coluna estreita ele ora cabia, ora não.
%%   Fio: fora da árvore (nenhum nó novo), sem conteúdo não marcado.
%%   Armadilhas achadas no caminho estão em comentário junto de cada
%%   macro: \rule num parágrafo com o tagging suspenso desbalanceia os
%%   hooks; o \vspace{-\parskip} da crpsp-base puxa o fio para cima do
%%   título, o multicols e o primeiro nome para dentro dele; o tipo `g'
%%   do ltcmd saiu do kernel.
%%
%% Compilar (2 passadas por variante):
%%   export TEXINPUTS="$HOME/repos/crpsp-latex/crpsp-book/desenvolvimento/v2:"
%%   for v in A B C E; do for i in 1 2; do lualatex -jobname=mwe_creditos_$v \
%%     "\def\variante{$v}%% mwe_creditos_titulo_invisivel.tex
%% Pergunta (P11 do plan-est_2026, segunda revisão de 24/09/2026):
%%   o modelo de créditos do Relatório de Gestão pode ir para a linha
%%   relatorio (PDF/UA-2) com a página de créditos SEM título impresso
%%   ("Créditos institucionais" etc.), mas com o título na árvore de
%%   tags, para o leitor de tela navegar?
%%
%% Variantes (por \def\variante na linha de comando):
%%   A  título invisível com texto real: \@startsection de pele zero,
%%      texto em Tr 3 (pdfrender), \smash\rlap, \topskip emulado.
%%   B  título sem conteúdo: \tagstructbegin{tag=H1,alt=...} vazio.
%%   C  sem título nenhum — a referência visual.
%% E, nas três, o que o modelo do RG pede além do título:
%%   - Plenário em duas colunas (o \parbox duplo do RG) pelo
%%     NomesDuasColunas do book;
%%   - \NomeNosCreditos e \CargoNosCreditos na interface do book 0.5.1,
%%     com o registro separado do nome (decisão de 17/09);
%%   - o fio sob o título de nível 2 (\CreditoInstitucionalUm) como
%%     artefato.
%%
%% RESULTADO (2026-09-24, TL2026, lualatex, veraPDF ua2):
%%   A  PASS. Árvore certa: Sect > H1 "Créditos institucionais" > Sect
%%      > H2 … ; o /T do H1 é o texto limpo; pdftotext extrai o título.
%%      A página sai IDÊNTICA à de C, pixel a pixel (pdftoppm -r 60 +
%%      diferença de imagem vazia), e o título seguinte cai em
%%      97,3956 pt nas duas. → É o padrão a portar.
%%   B  PASS no veraPDF, mas a estrutura quebra: o H1 fica solto, sem
%%      Sect, e o da 2ª página cai DENTRO do Sect de "Comissões
%%      permanentes" da 1ª. Um H1 sem conteúdo também depende de o
%%      leitor de tela ler /Alt em título, o que não é garantido.
%%      → Descartada.
%%   Nomes e cargos: interface do book 0.5.1 (book-crpsp_acessivel.sty
%%      :781-810) com a decisão de 17/09 aplicada — registro separado
%%      do nome, com o nome \NomeNosCreditos mantido (editor, 24/09):
%%        \NomeNosCreditos{nome}{registro}       forma nova
%%        \NomeNosCreditos{nome (CRP~06/nº)}     antiga, com aviso
%%        \CargoNosCreditos[registro]{nome}{cargo}
%%      Árvore: um P por pessoa, com o registro num Span próprio
%%      (14 Span no MWE: 10 nomes + 4 da Diretoria). A forma antiga sai
%%      como antes, nome e registro no mesmo MC.
%%   Duas colunas: NomesDuasColunas do book (multicols, \\ vira \par).
%%      Um P por pessoa, na ordem coluna 1 → coluna 2. O registro vai
%%      sempre para a linha de baixo, recuado (aprovado pelo editor em
%%      24/09). Não precisa do padrão de tagging do `paralelos'.
%%   Espaço entre pessoas: \crpspCreditosSep (0,3 \baselineskip) no
%%      CreditoInstitucional e no NomesDuasColunas. É exceção ao \parskip
%%      do WCAG numa lista de nomes — A CONFIRMAR COM O EDITOR.
%%      Do título ao primeiro nome: 5,67 pt nos dois ambientes.
%%   Porte (plan-est_2026, comum.tex §5, 24/09): acrescentou a nota
%%      opcional \NomeNosCreditos{nome}{registro}[licenciada] e a
%%      proibição de partir uma pessoa entre colunas/páginas (\clubpenalty,
%%      \widowpenalty e \interlinepenalty 10000 nos dois ambientes) —
%%      no pe26 o registro de uma pessoa foi parar na página seguinte.
%%   E  (terceira revisão, 24/09) cargos alinhados: \CargoNosCreditos
%%      com o nome numa \parbox de largura fixa (\crpspCreditosNome,
%%      0,56 \linewidth) e o cargo noutra, com o tagging do `paralelos'.
%%      PASS; por pessoa, Div > Div(P nome + Span registro) + Div(P
%%      cargo), na ordem certa, sem P vazio. O registro vai sempre na
%%      linha de baixo: na coluna estreita ele ora cabia, ora não.
%%   Fio: fora da árvore (nenhum nó novo), sem conteúdo não marcado.
%%   Armadilhas achadas no caminho estão em comentário junto de cada
%%   macro: \rule num parágrafo com o tagging suspenso desbalanceia os
%%   hooks; o \vspace{-\parskip} da crpsp-base puxa o fio para cima do
%%   título, o multicols e o primeiro nome para dentro dele; o tipo `g'
%%   do ltcmd saiu do kernel.
%%
%% Compilar (2 passadas por variante):
%%   export TEXINPUTS="$HOME/repos/crpsp-latex/crpsp-book/desenvolvimento/v2:"
%%   for v in A B C E; do for i in 1 2; do lualatex -jobname=mwe_creditos_$v \
%%     "\def\variante{$v}%% mwe_creditos_titulo_invisivel.tex
%% Pergunta (P11 do plan-est_2026, segunda revisão de 24/09/2026):
%%   o modelo de créditos do Relatório de Gestão pode ir para a linha
%%   relatorio (PDF/UA-2) com a página de créditos SEM título impresso
%%   ("Créditos institucionais" etc.), mas com o título na árvore de
%%   tags, para o leitor de tela navegar?
%%
%% Variantes (por \def\variante na linha de comando):
%%   A  título invisível com texto real: \@startsection de pele zero,
%%      texto em Tr 3 (pdfrender), \smash\rlap, \topskip emulado.
%%   B  título sem conteúdo: \tagstructbegin{tag=H1,alt=...} vazio.
%%   C  sem título nenhum — a referência visual.
%% E, nas três, o que o modelo do RG pede além do título:
%%   - Plenário em duas colunas (o \parbox duplo do RG) pelo
%%     NomesDuasColunas do book;
%%   - \NomeNosCreditos e \CargoNosCreditos na interface do book 0.5.1,
%%     com o registro separado do nome (decisão de 17/09);
%%   - o fio sob o título de nível 2 (\CreditoInstitucionalUm) como
%%     artefato.
%%
%% RESULTADO (2026-09-24, TL2026, lualatex, veraPDF ua2):
%%   A  PASS. Árvore certa: Sect > H1 "Créditos institucionais" > Sect
%%      > H2 … ; o /T do H1 é o texto limpo; pdftotext extrai o título.
%%      A página sai IDÊNTICA à de C, pixel a pixel (pdftoppm -r 60 +
%%      diferença de imagem vazia), e o título seguinte cai em
%%      97,3956 pt nas duas. → É o padrão a portar.
%%   B  PASS no veraPDF, mas a estrutura quebra: o H1 fica solto, sem
%%      Sect, e o da 2ª página cai DENTRO do Sect de "Comissões
%%      permanentes" da 1ª. Um H1 sem conteúdo também depende de o
%%      leitor de tela ler /Alt em título, o que não é garantido.
%%      → Descartada.
%%   Nomes e cargos: interface do book 0.5.1 (book-crpsp_acessivel.sty
%%      :781-810) com a decisão de 17/09 aplicada — registro separado
%%      do nome, com o nome \NomeNosCreditos mantido (editor, 24/09):
%%        \NomeNosCreditos{nome}{registro}       forma nova
%%        \NomeNosCreditos{nome (CRP~06/nº)}     antiga, com aviso
%%        \CargoNosCreditos[registro]{nome}{cargo}
%%      Árvore: um P por pessoa, com o registro num Span próprio
%%      (14 Span no MWE: 10 nomes + 4 da Diretoria). A forma antiga sai
%%      como antes, nome e registro no mesmo MC.
%%   Duas colunas: NomesDuasColunas do book (multicols, \\ vira \par).
%%      Um P por pessoa, na ordem coluna 1 → coluna 2. O registro vai
%%      sempre para a linha de baixo, recuado (aprovado pelo editor em
%%      24/09). Não precisa do padrão de tagging do `paralelos'.
%%   Espaço entre pessoas: \crpspCreditosSep (0,3 \baselineskip) no
%%      CreditoInstitucional e no NomesDuasColunas. É exceção ao \parskip
%%      do WCAG numa lista de nomes — A CONFIRMAR COM O EDITOR.
%%      Do título ao primeiro nome: 5,67 pt nos dois ambientes.
%%   Porte (plan-est_2026, comum.tex §5, 24/09): acrescentou a nota
%%      opcional \NomeNosCreditos{nome}{registro}[licenciada] e a
%%      proibição de partir uma pessoa entre colunas/páginas (\clubpenalty,
%%      \widowpenalty e \interlinepenalty 10000 nos dois ambientes) —
%%      no pe26 o registro de uma pessoa foi parar na página seguinte.
%%   E  (terceira revisão, 24/09) cargos alinhados: \CargoNosCreditos
%%      com o nome numa \parbox de largura fixa (\crpspCreditosNome,
%%      0,56 \linewidth) e o cargo noutra, com o tagging do `paralelos'.
%%      PASS; por pessoa, Div > Div(P nome + Span registro) + Div(P
%%      cargo), na ordem certa, sem P vazio. O registro vai sempre na
%%      linha de baixo: na coluna estreita ele ora cabia, ora não.
%%   Fio: fora da árvore (nenhum nó novo), sem conteúdo não marcado.
%%   Armadilhas achadas no caminho estão em comentário junto de cada
%%   macro: \rule num parágrafo com o tagging suspenso desbalanceia os
%%   hooks; o \vspace{-\parskip} da crpsp-base puxa o fio para cima do
%%   título, o multicols e o primeiro nome para dentro dele; o tipo `g'
%%   do ltcmd saiu do kernel.
%%
%% Compilar (2 passadas por variante):
%%   export TEXINPUTS="$HOME/repos/crpsp-latex/crpsp-book/desenvolvimento/v2:"
%%   for v in A B C E; do for i in 1 2; do lualatex -jobname=mwe_creditos_$v \
%%     "\def\variante{$v}\input{mwe_creditos_titulo_invisivel}"; done; done
%% Conferir:
%%   ~/verapdf/verapdf --format text --flavour ua2 mwe_creditos_A.pdf
%%   texlua $(kpsewhich show-pdf-tags.lua) --tree mwe_creditos_A.pdf
%%   pdftotext -bbox mwe_creditos_A.pdf - | grep -E '>(XVIII|Gerências)<'
%%
%% ⚠️ Aspas Unicode reais, nunca `` '': a linha compila com
%% \tracinglostchars=3 e a Atkinson não tem a crase.

\DocumentMetadata{
  lang        = pt-BR,
  pdfversion  = 2.0,
  pdfstandard = ua-2,
  tagging     = on,
}

\documentclass[externo]{crpsp-relatorio}

\providecommand{\variante}{A}

\usepackage{multicol}
\usepackage{pdfrender}

\makeatletter

% ── (A) título invisível com texto real ─────────────────────────
% O \@startsection da classe abre a estrutura de seção (section → H1,
% dentro de um Sect que engloba o que vem depois). O texto vai em Tr 3
% (pdfrender): não pinta, mas é conteúdo marcado — extração, cópia e
% leitor de tela o encontram.
%   - \smash\rlap: a linha do título tem altura, profundidade e
%     largura zero. Sem o \smash sobram 0,143 pt (a profundidade do
%     texto invisível).
%   - \texorpdfstring: sem ele o /T da estrutura sai
%     "TextRenderingMode=InvisibleCréditos institucionais".
%   - Depois do título, o próximo título (\creditosgrupo, \large) tem
%     de cair onde cairia no topo de uma página sem nada: linha de base
%     no \topskip. O \topskip só vale para a PRIMEIRA caixa da página,
%     e aqui a primeira é a linha invisível. A emulação:
%       \vskip-\topskip  volta a referência ao topo da mancha;
%       \prevdepth       faz a cola entre linhas valer
%                        \baselineskip(\large) − \prevdepth − h
%                        = \topskip − h, que é a cola do \topskip.
%     Medido (pdftotext -bbox): o título seguinte sai em 97,3956 pt,
%     igual à página sem título (variante C). Só com \vspace{-\parskip}
%     saía 3,7 pt acima; sem correção nenhuma, 17,9 pt abaixo.
%   ⚠️ A conta supõe que o próximo título é \large. Se o título de
%   nível 2 mudar de corpo, mudar aqui.
\newlength{\crpsp@bs@grupo}
\AtBeginDocument{{\large\global\crpsp@bs@grupo=\baselineskip}}
\newcommand{\crpsp@titulo@invisivel}[1]{%
  \textpdfrender{TextRenderingMode=Invisible}{#1}}
\newcommand{\creditosoculto}[1]{%
  \@startsection{section}{1}{0pt}{0pt}{1sp}%
    {\normalfont\fontsize{1pt}{0pt}\selectfont}%
    *{\texorpdfstring{\smash{\rlap{\crpsp@titulo@invisivel{\normalsize #1}}}}{#1}}%
  \par
  \vskip-\topskip
  \prevdepth=\dimexpr\crpsp@bs@grupo-\topskip\relax
}

% ── (B) título sem conteúdo, só /Alt ────────────────────────────
\newcommand{\creditosaltvazio}[1]{%
  \tagstructbegin{tag=H1,alt={#1}}\tagstructend}

% ── modelo do RG, com as cores e fontes da classe ────────────────
% \CreditoInstitucionalUm → nível 2 com fio; Dois → nível 3.
% (D) o fio é artefato (helper da classe, relatorio.sty §5).
\newcommand{\creditosgrupo}[1]{%
  \subsection*{#1}%
  % O \@afterheading da crpsp-base já deixou um \vspace{-\parskip} (o
  % parágrafo seguinte o devolveria). Sem devolver aqui, o fio sobe
  % 21,75 pt e sai ACIMA do título. Depois do fio o desconto é refeito,
  % para o que vier em seguida (lista, título) cair onde cairia.
  \vspace{\parskip}\vspace{2pt}%
  % em modo vertical (\hrule), não \rule num parágrafo: começar o
  % parágrafo com o tagging suspenso desbalanceia os hooks de
  % parágrafo ("begin (21) and end (24)"), um por fio.
  \crpsprel@artefatoinicio
  {\color{crpsp@cor3}\hrule height 1pt\relax}%
  \crpsprel@artefatofim
  \vspace{-\parskip}\vspace{0.4\baselineskip}%
}
% multicols logo depois de título: mesmo defeito do longtable
% (comum.tex §2 do plan-est_2026). Não abre parágrafo, não devolve o
% \parskip, e a primeira linha das colunas encavala no título.
\AddToHook{env/multicols/before}{\if@nobreak\vspace{\parskip}\fi}
\newcommand{\creditossub}[1]{\subsubsection*{#1}}

% ── nomes, cargos e registro: interface do book 0.5.1 ─────────────
% book-crpsp_acessivel.sty:781-810 (13/09/2026) tem \CargoNosCreditos
% de dois argumentos, \NomeNosCreditos de um e o NomesDuasColunas. A
% decisão de 17/09 (briefings/funcionalidade-creditos-institucionais.md)
% separa o registro do nome, para que ele não seja lido como parte do
% nome. Aqui o registro vira uma estrutura própria (Span) dentro do
% parágrafo da pessoa.
%
% O nome fica \NomeNosCreditos (decisão do editor, 24/09/2026), e a
% forma antiga continua aceita pelo mesmo nome: ver adiante.
\ExplSyntaxOn
% Fecha o MC do parágrafo, abre Span com MC próprio, e reabre o MC do
% parágrafo: o par push/pop do tagpdf existe para isto.
\cs_new_protected:Npn \crpsp_creditos_registro:n #1
  {
    \tag_mc_end_push:
    \tag_struct_begin:n { tag = Span }
    \tag_mc_begin:n { }
    #1
    \tag_mc_end:
    \tag_struct_end:
    \tag_mc_begin_pop:n { }
  }
\cs_new_eq:NN \crpsp@registro \crpsp_creditos_registro:n
\ExplSyntaxOff
% Separador entre nome e registro: espaço na largura da mancha; nas
% duas colunas, quebra de linha com recuo (aprovado pelo editor em
% 24/09: registro sempre embaixo, nunca solto pela largura da coluna).
\newcommand{\crpsp@registro@sep}{ }
\newcommand{\crpsp@registro@fmt}[1]{(CRP~#1)}

% \NomeNosCreditos{nome}{registro}: uma pessoa, um parágrafo.
%   \NomeNosCreditos{Beatris Guarita Dotta}{06/143345}
% Forma antiga (book 0.5.1), um argumento com o registro no texto —
% deprecada, com aviso:
%   \NomeNosCreditos{Beatris Guarita Dotta (CRP~06/143345)}
% As duas convivem no mesmo nome: depois do primeiro argumento, se o
% próximo token (pulando espaços) for `{', é a forma nova. O tipo `g'
% do ltcmd faria o mesmo, mas saiu do kernel ("Invalid argument type
% 'g'") e só volta com o xparse obsoleto. ⚠️ A armadilha é a mesma do
% `g': na forma antiga, se o que vem depois começar por `{', é engolido
% como registro. Nos .tex existentes o que vem depois é outro
% \NomeNosCreditos, \\, `}' ou linha em branco. Quando a forma antiga
% sair, trocar por `m m'.
\ExplSyntaxOn
\NewDocumentCommand{\NomeNosCreditos}{m}
  {
    \peek_catcode_ignore_spaces:NTF \c_group_begin_token
      { \crpsp_creditos_nome:nn {#1} }
      {
        \msg_warning:nn { crpsp-creditos } { nome-antigo }
        \crpsp_creditos_nome_antigo:n {#1}
      }
  }
\cs_new_protected:Npn \crpsp_creditos_nome:nn #1#2
  {
    \par\noindent\hangindent=1em\hangafter=1\relax
    #1 \crpsp@registro@sep \crpsp@registro { \crpsp@registro@fmt {#2} }
    \par
  }
\cs_new_protected:Npn \crpsp_creditos_nome_antigo:n #1
  { \par\noindent #1 \par }
\msg_new:nnn { crpsp-creditos } { nome-antigo }
  {
    \iow_char:N\\NomeNosCreditos~de~um~argumento~está~deprecado:~use~
    \iow_char:N\\NomeNosCreditos\{nome\}\{registro\}.
  }
\ExplSyntaxOff

% \CargoNosCreditos[registro]{nome}{cargo}: os dois obrigatórios são
% os do book 0.5.1; o registro entra como opcional na frente para não
% quebrar os .tex que já usam a forma de dois. Visual do RG
% (relatorio-gestao.sty:246): nome em destaque, cargo ao lado. As duas
% \parbox 0,6/0,4 do book não vieram: no RG o cargo segue o nome.
\NewDocumentCommand{\CargoNosCreditos}{o m m}{%
  \par\noindent\hangindent=1em\hangafter=1\relax
  \textbf{#2}%
  \IfValueT{#1}{ \crpsp@registro{\crpsp@registro@fmt{#1}}}%
  \quad #3\par}

% ── (E) cargos alinhados: \CargoNosCreditos com coluna fixa ───────
% Terceira revisão do plan-est_2026, item 1: o book 0.5.1 tinha duas
% \parbox (0,6/0,4) e os cargos saíam alinhados na vertical. Aqui a
% mesma forma, com o tagging do `paralelos' (relatorio.sty §6): a
% pessoa é um Div; cada \parbox abre o seu parágrafo (paraon dentro,
% \par e paraoff antes de fechar). Leitura: nome → registro → cargo.
%   \crpspCreditosNome   largura da coluna do nome (com o registro)
%   \crpspCreditosVao    espaço entre as colunas
\newlength{\crpspCreditosNome}
\newlength{\crpspCreditosVao}
\AtBeginDocument{%
  \setlength{\crpspCreditosNome}{0.56\linewidth}%
  \setlength{\crpspCreditosVao}{1em}}
\newcommand{\CargoAlinhado}[3]{% #1 registro (ou vazio), #2 nome, #3 cargo
  \par
  \crpsprel@paraoff
  \crpsprel@tagbegin{Div}%
  \noindent
  \parbox[t]{\crpspCreditosNome}{%
    \raggedright\hangindent=1em\hangafter=1\relax
    \crpsprel@paraon
    \textbf{#2}%
    % registro sempre na linha de baixo, como nas duas colunas: na
    % coluna estreita ele ora cabia, ora não, e a lista ficava irregular
    \ifx\relax#1\relax\else\newline\crpsp@registro{\crpsp@registro@fmt{#1}}\fi
    \par\crpsprel@paraoff}%
  \hspace{\crpspCreditosVao}%
  \parbox[t]{\dimexpr\linewidth-\crpspCreditosNome-\crpspCreditosVao\relax}{%
    \raggedright
    \crpsprel@paraon
    #3\par\crpsprel@paraoff}%
  \par
  \crpsprel@tagend
  \crpsprel@paraon
}
\if E\variante
  \RenewDocumentCommand{\CargoNosCreditos}{o m m}{%
    \IfValueTF{#1}{\CargoAlinhado{#1}}{\CargoAlinhado{}}{#2}{#3}}
\fi

% CreditoInstitucional: o ambiente do modelo (RG e book) que envolve
% um bloco de créditos. Cada pessoa é um parágrafo, e o \parskip do
% WCAG (1,5 × a entrelinha, 21,75 pt) abria a lista demais. Aqui ele
% cai para \crpspCreditosSep. Como no sumário (comum.tex §4 do
% plan-est_2026), é exceção ao WCAG 1.4.8 numa lista de nomes, não em
% texto corrido: CONFIRMAR COM O EDITOR.
\newlength{\crpspCreditosSep}
\setlength{\crpspCreditosSep}{0.3\baselineskip}
% ⚠️ Logo depois de título, o \@afterheading da crpsp-base já
% descontou o \parskip INTEIRO, e o primeiro parágrafo daqui só devolve
% o \crpspCreditosSep: o primeiro nome encavalava no título. A diferença
% é devolvida na abertura.
\newenvironment{CreditoInstitucional}
  {\par
   \if@nobreak\vspace{\dimexpr\parskip-\crpspCreditosSep\relax}\fi
   \begingroup\setlength{\parskip}{\crpspCreditosSep}}
  {\par\endgroup}

% NomesDuasColunas: o do book (multicols, \\ vira \par), com o
% separador trocado por quebra de linha.
% O \multicolsep (acima e abaixo das colunas) punha 17,6 pt entre o
% título e o primeiro nome, contra 5,7 pt na Diretoria. Aqui ele fica
% zerado: o espaço de cima é o do hook env/multicols/before e o de
% baixo é o do título seguinte.
\newenvironment{NomesDuasColunas}{%
  \par\begingroup
  \setlength{\multicolsep}{0pt}%
  \begin{multicols}{2}%
  \begingroup
  \raggedright
  \setlength{\parskip}{\crpspCreditosSep}%
  \let\\\par
  \renewcommand{\crpsp@registro@sep}{\newline}%
}{%
  \endgroup
  \end{multicols}%
  \endgroup
}

\makeatother

\hypersetup{pdftitle={MWE — créditos com título invisível}}

\begin{document}

% ════════════════════════ página de créditos
\if B\variante\else\if C\variante\else\creditosoculto{Créditos institucionais}\fi\fi
\if B\variante\creditosaltvazio{Créditos institucionais}\fi

\creditosgrupo{XVIII Plenário (2025–2028)}

\creditossub{Diretoria}
\begin{CreditoInstitucional}
\CargoNosCreditos[06/31093]{Valéria Campinas Braunstein}{presidenta}
\CargoNosCreditos[06/149691]{Victória Soares Vidal}{vice-presidenta}
\CargoNosCreditos[06/148611]{Fabiana Macena Luiz}{secretária}
\CargoNosCreditos[06/159297]{Genildo Gomes de Sousa}{tesoureiro}
\end{CreditoInstitucional}

\creditossub{Conselheiras/os efetivas/os}
\begin{NomesDuasColunas}
  \NomeNosCreditos{Cláudia Cristina Lofrano}{06/44926}
  \NomeNosCreditos{Débora Nascimento Santos}{06/144738}
  \NomeNosCreditos{Fabiana Macena Luiz}{06/148611}
  \NomeNosCreditos{Gabriel Basílio Barbosa Costa}{06/185699}
  \NomeNosCreditos{Genildo Gomes de Sousa}{06/159297}
  \NomeNosCreditos{Luísa Thomazini de Freitas}{06/159754}
  \NomeNosCreditos{Luke Ribeiro Mazzei França Barros}{06/188231}
  \NomeNosCreditos{Marilia Capponi}{06/81224}
  \NomeNosCreditos{Natali de Souza Nascimento}{06/134212}
  \NomeNosCreditos{Paula Andréia de Carvalho Jonas}{06/62340}
\end{NomesDuasColunas}

\creditossub{Conselheiras/os suplentes (forma antiga, deprecada)}
\begin{NomesDuasColunas}
  \NomeNosCreditos{Beatris Guarita Dotta (CRP~06/143345)}
  \NomeNosCreditos{Bruna Pessenda (CRP~06/137732)}
\end{NomesDuasColunas}

\creditosgrupo{Comissões permanentes}

\creditossub{Comissão de Ética (COE)}
\begin{CreditoInstitucional}
\CargoNosCreditos{Cláudia Cristina Lofrano}{presidenta}
\end{CreditoInstitucional}

\clearpage
% ════════════════════════ segunda página: título oculto no topo,
% seguido direto de texto, para medir que nada foi empurrado
\if B\variante\else\if C\variante\else\creditosoculto{Créditos técnico-administrativos}\fi\fi
\if B\variante\creditosaltvazio{Créditos técnico-administrativos}\fi
\creditosgrupo{Gerências}
\begin{CreditoInstitucional}
\CargoNosCreditos{Vanessa Valente}{gerente de Administração e Tecnologia da Informação}
\CargoNosCreditos{Lilian Suzuki}{gerente Técnica e Política (interina)}
\CargoNosCreditos[06/110639]{Maria Carolina Vitalino Pinto Ferraz Cabau}{subcoordenadora}
\CargoNosCreditos{Adolfo Barros Benevenuto}{coordenador de Tecnologia da Informação e Comunicação}
\end{CreditoInstitucional}

\end{document}
"; done; done
%% Conferir:
%%   ~/verapdf/verapdf --format text --flavour ua2 mwe_creditos_A.pdf
%%   texlua $(kpsewhich show-pdf-tags.lua) --tree mwe_creditos_A.pdf
%%   pdftotext -bbox mwe_creditos_A.pdf - | grep -E '>(XVIII|Gerências)<'
%%
%% ⚠️ Aspas Unicode reais, nunca `` '': a linha compila com
%% \tracinglostchars=3 e a Atkinson não tem a crase.

\DocumentMetadata{
  lang        = pt-BR,
  pdfversion  = 2.0,
  pdfstandard = ua-2,
  tagging     = on,
}

\documentclass[externo]{crpsp-relatorio}

\providecommand{\variante}{A}

\usepackage{multicol}
\usepackage{pdfrender}

\makeatletter

% ── (A) título invisível com texto real ─────────────────────────
% O \@startsection da classe abre a estrutura de seção (section → H1,
% dentro de um Sect que engloba o que vem depois). O texto vai em Tr 3
% (pdfrender): não pinta, mas é conteúdo marcado — extração, cópia e
% leitor de tela o encontram.
%   - \smash\rlap: a linha do título tem altura, profundidade e
%     largura zero. Sem o \smash sobram 0,143 pt (a profundidade do
%     texto invisível).
%   - \texorpdfstring: sem ele o /T da estrutura sai
%     "TextRenderingMode=InvisibleCréditos institucionais".
%   - Depois do título, o próximo título (\creditosgrupo, \large) tem
%     de cair onde cairia no topo de uma página sem nada: linha de base
%     no \topskip. O \topskip só vale para a PRIMEIRA caixa da página,
%     e aqui a primeira é a linha invisível. A emulação:
%       \vskip-\topskip  volta a referência ao topo da mancha;
%       \prevdepth       faz a cola entre linhas valer
%                        \baselineskip(\large) − \prevdepth − h
%                        = \topskip − h, que é a cola do \topskip.
%     Medido (pdftotext -bbox): o título seguinte sai em 97,3956 pt,
%     igual à página sem título (variante C). Só com \vspace{-\parskip}
%     saía 3,7 pt acima; sem correção nenhuma, 17,9 pt abaixo.
%   ⚠️ A conta supõe que o próximo título é \large. Se o título de
%   nível 2 mudar de corpo, mudar aqui.
\newlength{\crpsp@bs@grupo}
\AtBeginDocument{{\large\global\crpsp@bs@grupo=\baselineskip}}
\newcommand{\crpsp@titulo@invisivel}[1]{%
  \textpdfrender{TextRenderingMode=Invisible}{#1}}
\newcommand{\creditosoculto}[1]{%
  \@startsection{section}{1}{0pt}{0pt}{1sp}%
    {\normalfont\fontsize{1pt}{0pt}\selectfont}%
    *{\texorpdfstring{\smash{\rlap{\crpsp@titulo@invisivel{\normalsize #1}}}}{#1}}%
  \par
  \vskip-\topskip
  \prevdepth=\dimexpr\crpsp@bs@grupo-\topskip\relax
}

% ── (B) título sem conteúdo, só /Alt ────────────────────────────
\newcommand{\creditosaltvazio}[1]{%
  \tagstructbegin{tag=H1,alt={#1}}\tagstructend}

% ── modelo do RG, com as cores e fontes da classe ────────────────
% \CreditoInstitucionalUm → nível 2 com fio; Dois → nível 3.
% (D) o fio é artefato (helper da classe, relatorio.sty §5).
\newcommand{\creditosgrupo}[1]{%
  \subsection*{#1}%
  % O \@afterheading da crpsp-base já deixou um \vspace{-\parskip} (o
  % parágrafo seguinte o devolveria). Sem devolver aqui, o fio sobe
  % 21,75 pt e sai ACIMA do título. Depois do fio o desconto é refeito,
  % para o que vier em seguida (lista, título) cair onde cairia.
  \vspace{\parskip}\vspace{2pt}%
  % em modo vertical (\hrule), não \rule num parágrafo: começar o
  % parágrafo com o tagging suspenso desbalanceia os hooks de
  % parágrafo ("begin (21) and end (24)"), um por fio.
  \crpsprel@artefatoinicio
  {\color{crpsp@cor3}\hrule height 1pt\relax}%
  \crpsprel@artefatofim
  \vspace{-\parskip}\vspace{0.4\baselineskip}%
}
% multicols logo depois de título: mesmo defeito do longtable
% (comum.tex §2 do plan-est_2026). Não abre parágrafo, não devolve o
% \parskip, e a primeira linha das colunas encavala no título.
\AddToHook{env/multicols/before}{\if@nobreak\vspace{\parskip}\fi}
\newcommand{\creditossub}[1]{\subsubsection*{#1}}

% ── nomes, cargos e registro: interface do book 0.5.1 ─────────────
% book-crpsp_acessivel.sty:781-810 (13/09/2026) tem \CargoNosCreditos
% de dois argumentos, \NomeNosCreditos de um e o NomesDuasColunas. A
% decisão de 17/09 (briefings/funcionalidade-creditos-institucionais.md)
% separa o registro do nome, para que ele não seja lido como parte do
% nome. Aqui o registro vira uma estrutura própria (Span) dentro do
% parágrafo da pessoa.
%
% O nome fica \NomeNosCreditos (decisão do editor, 24/09/2026), e a
% forma antiga continua aceita pelo mesmo nome: ver adiante.
\ExplSyntaxOn
% Fecha o MC do parágrafo, abre Span com MC próprio, e reabre o MC do
% parágrafo: o par push/pop do tagpdf existe para isto.
\cs_new_protected:Npn \crpsp_creditos_registro:n #1
  {
    \tag_mc_end_push:
    \tag_struct_begin:n { tag = Span }
    \tag_mc_begin:n { }
    #1
    \tag_mc_end:
    \tag_struct_end:
    \tag_mc_begin_pop:n { }
  }
\cs_new_eq:NN \crpsp@registro \crpsp_creditos_registro:n
\ExplSyntaxOff
% Separador entre nome e registro: espaço na largura da mancha; nas
% duas colunas, quebra de linha com recuo (aprovado pelo editor em
% 24/09: registro sempre embaixo, nunca solto pela largura da coluna).
\newcommand{\crpsp@registro@sep}{ }
\newcommand{\crpsp@registro@fmt}[1]{(CRP~#1)}

% \NomeNosCreditos{nome}{registro}: uma pessoa, um parágrafo.
%   \NomeNosCreditos{Beatris Guarita Dotta}{06/143345}
% Forma antiga (book 0.5.1), um argumento com o registro no texto —
% deprecada, com aviso:
%   \NomeNosCreditos{Beatris Guarita Dotta (CRP~06/143345)}
% As duas convivem no mesmo nome: depois do primeiro argumento, se o
% próximo token (pulando espaços) for `{', é a forma nova. O tipo `g'
% do ltcmd faria o mesmo, mas saiu do kernel ("Invalid argument type
% 'g'") e só volta com o xparse obsoleto. ⚠️ A armadilha é a mesma do
% `g': na forma antiga, se o que vem depois começar por `{', é engolido
% como registro. Nos .tex existentes o que vem depois é outro
% \NomeNosCreditos, \\, `}' ou linha em branco. Quando a forma antiga
% sair, trocar por `m m'.
\ExplSyntaxOn
\NewDocumentCommand{\NomeNosCreditos}{m}
  {
    \peek_catcode_ignore_spaces:NTF \c_group_begin_token
      { \crpsp_creditos_nome:nn {#1} }
      {
        \msg_warning:nn { crpsp-creditos } { nome-antigo }
        \crpsp_creditos_nome_antigo:n {#1}
      }
  }
\cs_new_protected:Npn \crpsp_creditos_nome:nn #1#2
  {
    \par\noindent\hangindent=1em\hangafter=1\relax
    #1 \crpsp@registro@sep \crpsp@registro { \crpsp@registro@fmt {#2} }
    \par
  }
\cs_new_protected:Npn \crpsp_creditos_nome_antigo:n #1
  { \par\noindent #1 \par }
\msg_new:nnn { crpsp-creditos } { nome-antigo }
  {
    \iow_char:N\\NomeNosCreditos~de~um~argumento~está~deprecado:~use~
    \iow_char:N\\NomeNosCreditos\{nome\}\{registro\}.
  }
\ExplSyntaxOff

% \CargoNosCreditos[registro]{nome}{cargo}: os dois obrigatórios são
% os do book 0.5.1; o registro entra como opcional na frente para não
% quebrar os .tex que já usam a forma de dois. Visual do RG
% (relatorio-gestao.sty:246): nome em destaque, cargo ao lado. As duas
% \parbox 0,6/0,4 do book não vieram: no RG o cargo segue o nome.
\NewDocumentCommand{\CargoNosCreditos}{o m m}{%
  \par\noindent\hangindent=1em\hangafter=1\relax
  \textbf{#2}%
  \IfValueT{#1}{ \crpsp@registro{\crpsp@registro@fmt{#1}}}%
  \quad #3\par}

% ── (E) cargos alinhados: \CargoNosCreditos com coluna fixa ───────
% Terceira revisão do plan-est_2026, item 1: o book 0.5.1 tinha duas
% \parbox (0,6/0,4) e os cargos saíam alinhados na vertical. Aqui a
% mesma forma, com o tagging do `paralelos' (relatorio.sty §6): a
% pessoa é um Div; cada \parbox abre o seu parágrafo (paraon dentro,
% \par e paraoff antes de fechar). Leitura: nome → registro → cargo.
%   \crpspCreditosNome   largura da coluna do nome (com o registro)
%   \crpspCreditosVao    espaço entre as colunas
\newlength{\crpspCreditosNome}
\newlength{\crpspCreditosVao}
\AtBeginDocument{%
  \setlength{\crpspCreditosNome}{0.56\linewidth}%
  \setlength{\crpspCreditosVao}{1em}}
\newcommand{\CargoAlinhado}[3]{% #1 registro (ou vazio), #2 nome, #3 cargo
  \par
  \crpsprel@paraoff
  \crpsprel@tagbegin{Div}%
  \noindent
  \parbox[t]{\crpspCreditosNome}{%
    \raggedright\hangindent=1em\hangafter=1\relax
    \crpsprel@paraon
    \textbf{#2}%
    % registro sempre na linha de baixo, como nas duas colunas: na
    % coluna estreita ele ora cabia, ora não, e a lista ficava irregular
    \ifx\relax#1\relax\else\newline\crpsp@registro{\crpsp@registro@fmt{#1}}\fi
    \par\crpsprel@paraoff}%
  \hspace{\crpspCreditosVao}%
  \parbox[t]{\dimexpr\linewidth-\crpspCreditosNome-\crpspCreditosVao\relax}{%
    \raggedright
    \crpsprel@paraon
    #3\par\crpsprel@paraoff}%
  \par
  \crpsprel@tagend
  \crpsprel@paraon
}
\if E\variante
  \RenewDocumentCommand{\CargoNosCreditos}{o m m}{%
    \IfValueTF{#1}{\CargoAlinhado{#1}}{\CargoAlinhado{}}{#2}{#3}}
\fi

% CreditoInstitucional: o ambiente do modelo (RG e book) que envolve
% um bloco de créditos. Cada pessoa é um parágrafo, e o \parskip do
% WCAG (1,5 × a entrelinha, 21,75 pt) abria a lista demais. Aqui ele
% cai para \crpspCreditosSep. Como no sumário (comum.tex §4 do
% plan-est_2026), é exceção ao WCAG 1.4.8 numa lista de nomes, não em
% texto corrido: CONFIRMAR COM O EDITOR.
\newlength{\crpspCreditosSep}
\setlength{\crpspCreditosSep}{0.3\baselineskip}
% ⚠️ Logo depois de título, o \@afterheading da crpsp-base já
% descontou o \parskip INTEIRO, e o primeiro parágrafo daqui só devolve
% o \crpspCreditosSep: o primeiro nome encavalava no título. A diferença
% é devolvida na abertura.
\newenvironment{CreditoInstitucional}
  {\par
   \if@nobreak\vspace{\dimexpr\parskip-\crpspCreditosSep\relax}\fi
   \begingroup\setlength{\parskip}{\crpspCreditosSep}}
  {\par\endgroup}

% NomesDuasColunas: o do book (multicols, \\ vira \par), com o
% separador trocado por quebra de linha.
% O \multicolsep (acima e abaixo das colunas) punha 17,6 pt entre o
% título e o primeiro nome, contra 5,7 pt na Diretoria. Aqui ele fica
% zerado: o espaço de cima é o do hook env/multicols/before e o de
% baixo é o do título seguinte.
\newenvironment{NomesDuasColunas}{%
  \par\begingroup
  \setlength{\multicolsep}{0pt}%
  \begin{multicols}{2}%
  \begingroup
  \raggedright
  \setlength{\parskip}{\crpspCreditosSep}%
  \let\\\par
  \renewcommand{\crpsp@registro@sep}{\newline}%
}{%
  \endgroup
  \end{multicols}%
  \endgroup
}

\makeatother

\hypersetup{pdftitle={MWE — créditos com título invisível}}

\begin{document}

% ════════════════════════ página de créditos
\if B\variante\else\if C\variante\else\creditosoculto{Créditos institucionais}\fi\fi
\if B\variante\creditosaltvazio{Créditos institucionais}\fi

\creditosgrupo{XVIII Plenário (2025–2028)}

\creditossub{Diretoria}
\begin{CreditoInstitucional}
\CargoNosCreditos[06/31093]{Valéria Campinas Braunstein}{presidenta}
\CargoNosCreditos[06/149691]{Victória Soares Vidal}{vice-presidenta}
\CargoNosCreditos[06/148611]{Fabiana Macena Luiz}{secretária}
\CargoNosCreditos[06/159297]{Genildo Gomes de Sousa}{tesoureiro}
\end{CreditoInstitucional}

\creditossub{Conselheiras/os efetivas/os}
\begin{NomesDuasColunas}
  \NomeNosCreditos{Cláudia Cristina Lofrano}{06/44926}
  \NomeNosCreditos{Débora Nascimento Santos}{06/144738}
  \NomeNosCreditos{Fabiana Macena Luiz}{06/148611}
  \NomeNosCreditos{Gabriel Basílio Barbosa Costa}{06/185699}
  \NomeNosCreditos{Genildo Gomes de Sousa}{06/159297}
  \NomeNosCreditos{Luísa Thomazini de Freitas}{06/159754}
  \NomeNosCreditos{Luke Ribeiro Mazzei França Barros}{06/188231}
  \NomeNosCreditos{Marilia Capponi}{06/81224}
  \NomeNosCreditos{Natali de Souza Nascimento}{06/134212}
  \NomeNosCreditos{Paula Andréia de Carvalho Jonas}{06/62340}
\end{NomesDuasColunas}

\creditossub{Conselheiras/os suplentes (forma antiga, deprecada)}
\begin{NomesDuasColunas}
  \NomeNosCreditos{Beatris Guarita Dotta (CRP~06/143345)}
  \NomeNosCreditos{Bruna Pessenda (CRP~06/137732)}
\end{NomesDuasColunas}

\creditosgrupo{Comissões permanentes}

\creditossub{Comissão de Ética (COE)}
\begin{CreditoInstitucional}
\CargoNosCreditos{Cláudia Cristina Lofrano}{presidenta}
\end{CreditoInstitucional}

\clearpage
% ════════════════════════ segunda página: título oculto no topo,
% seguido direto de texto, para medir que nada foi empurrado
\if B\variante\else\if C\variante\else\creditosoculto{Créditos técnico-administrativos}\fi\fi
\if B\variante\creditosaltvazio{Créditos técnico-administrativos}\fi
\creditosgrupo{Gerências}
\begin{CreditoInstitucional}
\CargoNosCreditos{Vanessa Valente}{gerente de Administração e Tecnologia da Informação}
\CargoNosCreditos{Lilian Suzuki}{gerente Técnica e Política (interina)}
\CargoNosCreditos[06/110639]{Maria Carolina Vitalino Pinto Ferraz Cabau}{subcoordenadora}
\CargoNosCreditos{Adolfo Barros Benevenuto}{coordenador de Tecnologia da Informação e Comunicação}
\end{CreditoInstitucional}

\end{document}
"; done; done
%% Conferir:
%%   ~/verapdf/verapdf --format text --flavour ua2 mwe_creditos_A.pdf
%%   texlua $(kpsewhich show-pdf-tags.lua) --tree mwe_creditos_A.pdf
%%   pdftotext -bbox mwe_creditos_A.pdf - | grep -E '>(XVIII|Gerências)<'
%%
%% ⚠️ Aspas Unicode reais, nunca `` '': a linha compila com
%% \tracinglostchars=3 e a Atkinson não tem a crase.

\DocumentMetadata{
  lang        = pt-BR,
  pdfversion  = 2.0,
  pdfstandard = ua-2,
  tagging     = on,
}

\documentclass[externo]{crpsp-relatorio}

\providecommand{\variante}{A}

\usepackage{multicol}
\usepackage{pdfrender}

\makeatletter

% ── (A) título invisível com texto real ─────────────────────────
% O \@startsection da classe abre a estrutura de seção (section → H1,
% dentro de um Sect que engloba o que vem depois). O texto vai em Tr 3
% (pdfrender): não pinta, mas é conteúdo marcado — extração, cópia e
% leitor de tela o encontram.
%   - \smash\rlap: a linha do título tem altura, profundidade e
%     largura zero. Sem o \smash sobram 0,143 pt (a profundidade do
%     texto invisível).
%   - \texorpdfstring: sem ele o /T da estrutura sai
%     "TextRenderingMode=InvisibleCréditos institucionais".
%   - Depois do título, o próximo título (\creditosgrupo, \large) tem
%     de cair onde cairia no topo de uma página sem nada: linha de base
%     no \topskip. O \topskip só vale para a PRIMEIRA caixa da página,
%     e aqui a primeira é a linha invisível. A emulação:
%       \vskip-\topskip  volta a referência ao topo da mancha;
%       \prevdepth       faz a cola entre linhas valer
%                        \baselineskip(\large) − \prevdepth − h
%                        = \topskip − h, que é a cola do \topskip.
%     Medido (pdftotext -bbox): o título seguinte sai em 97,3956 pt,
%     igual à página sem título (variante C). Só com \vspace{-\parskip}
%     saía 3,7 pt acima; sem correção nenhuma, 17,9 pt abaixo.
%   ⚠️ A conta supõe que o próximo título é \large. Se o título de
%   nível 2 mudar de corpo, mudar aqui.
\newlength{\crpsp@bs@grupo}
\AtBeginDocument{{\large\global\crpsp@bs@grupo=\baselineskip}}
\newcommand{\crpsp@titulo@invisivel}[1]{%
  \textpdfrender{TextRenderingMode=Invisible}{#1}}
\newcommand{\creditosoculto}[1]{%
  \@startsection{section}{1}{0pt}{0pt}{1sp}%
    {\normalfont\fontsize{1pt}{0pt}\selectfont}%
    *{\texorpdfstring{\smash{\rlap{\crpsp@titulo@invisivel{\normalsize #1}}}}{#1}}%
  \par
  \vskip-\topskip
  \prevdepth=\dimexpr\crpsp@bs@grupo-\topskip\relax
}

% ── (B) título sem conteúdo, só /Alt ────────────────────────────
\newcommand{\creditosaltvazio}[1]{%
  \tagstructbegin{tag=H1,alt={#1}}\tagstructend}

% ── modelo do RG, com as cores e fontes da classe ────────────────
% \CreditoInstitucionalUm → nível 2 com fio; Dois → nível 3.
% (D) o fio é artefato (helper da classe, relatorio.sty §5).
\newcommand{\creditosgrupo}[1]{%
  \subsection*{#1}%
  % O \@afterheading da crpsp-base já deixou um \vspace{-\parskip} (o
  % parágrafo seguinte o devolveria). Sem devolver aqui, o fio sobe
  % 21,75 pt e sai ACIMA do título. Depois do fio o desconto é refeito,
  % para o que vier em seguida (lista, título) cair onde cairia.
  \vspace{\parskip}\vspace{2pt}%
  % em modo vertical (\hrule), não \rule num parágrafo: começar o
  % parágrafo com o tagging suspenso desbalanceia os hooks de
  % parágrafo ("begin (21) and end (24)"), um por fio.
  \crpsprel@artefatoinicio
  {\color{crpsp@cor3}\hrule height 1pt\relax}%
  \crpsprel@artefatofim
  \vspace{-\parskip}\vspace{0.4\baselineskip}%
}
% multicols logo depois de título: mesmo defeito do longtable
% (comum.tex §2 do plan-est_2026). Não abre parágrafo, não devolve o
% \parskip, e a primeira linha das colunas encavala no título.
\AddToHook{env/multicols/before}{\if@nobreak\vspace{\parskip}\fi}
\newcommand{\creditossub}[1]{\subsubsection*{#1}}

% ── nomes, cargos e registro: interface do book 0.5.1 ─────────────
% book-crpsp_acessivel.sty:781-810 (13/09/2026) tem \CargoNosCreditos
% de dois argumentos, \NomeNosCreditos de um e o NomesDuasColunas. A
% decisão de 17/09 (briefings/funcionalidade-creditos-institucionais.md)
% separa o registro do nome, para que ele não seja lido como parte do
% nome. Aqui o registro vira uma estrutura própria (Span) dentro do
% parágrafo da pessoa.
%
% O nome fica \NomeNosCreditos (decisão do editor, 24/09/2026), e a
% forma antiga continua aceita pelo mesmo nome: ver adiante.
\ExplSyntaxOn
% Fecha o MC do parágrafo, abre Span com MC próprio, e reabre o MC do
% parágrafo: o par push/pop do tagpdf existe para isto.
\cs_new_protected:Npn \crpsp_creditos_registro:n #1
  {
    \tag_mc_end_push:
    \tag_struct_begin:n { tag = Span }
    \tag_mc_begin:n { }
    #1
    \tag_mc_end:
    \tag_struct_end:
    \tag_mc_begin_pop:n { }
  }
\cs_new_eq:NN \crpsp@registro \crpsp_creditos_registro:n
\ExplSyntaxOff
% Separador entre nome e registro: espaço na largura da mancha; nas
% duas colunas, quebra de linha com recuo (aprovado pelo editor em
% 24/09: registro sempre embaixo, nunca solto pela largura da coluna).
\newcommand{\crpsp@registro@sep}{ }
\newcommand{\crpsp@registro@fmt}[1]{(CRP~#1)}

% \NomeNosCreditos{nome}{registro}: uma pessoa, um parágrafo.
%   \NomeNosCreditos{Beatris Guarita Dotta}{06/143345}
% Forma antiga (book 0.5.1), um argumento com o registro no texto —
% deprecada, com aviso:
%   \NomeNosCreditos{Beatris Guarita Dotta (CRP~06/143345)}
% As duas convivem no mesmo nome: depois do primeiro argumento, se o
% próximo token (pulando espaços) for `{', é a forma nova. O tipo `g'
% do ltcmd faria o mesmo, mas saiu do kernel ("Invalid argument type
% 'g'") e só volta com o xparse obsoleto. ⚠️ A armadilha é a mesma do
% `g': na forma antiga, se o que vem depois começar por `{', é engolido
% como registro. Nos .tex existentes o que vem depois é outro
% \NomeNosCreditos, \\, `}' ou linha em branco. Quando a forma antiga
% sair, trocar por `m m'.
\ExplSyntaxOn
\NewDocumentCommand{\NomeNosCreditos}{m}
  {
    \peek_catcode_ignore_spaces:NTF \c_group_begin_token
      { \crpsp_creditos_nome:nn {#1} }
      {
        \msg_warning:nn { crpsp-creditos } { nome-antigo }
        \crpsp_creditos_nome_antigo:n {#1}
      }
  }
\cs_new_protected:Npn \crpsp_creditos_nome:nn #1#2
  {
    \par\noindent\hangindent=1em\hangafter=1\relax
    #1 \crpsp@registro@sep \crpsp@registro { \crpsp@registro@fmt {#2} }
    \par
  }
\cs_new_protected:Npn \crpsp_creditos_nome_antigo:n #1
  { \par\noindent #1 \par }
\msg_new:nnn { crpsp-creditos } { nome-antigo }
  {
    \iow_char:N\\NomeNosCreditos~de~um~argumento~está~deprecado:~use~
    \iow_char:N\\NomeNosCreditos\{nome\}\{registro\}.
  }
\ExplSyntaxOff

% \CargoNosCreditos[registro]{nome}{cargo}: os dois obrigatórios são
% os do book 0.5.1; o registro entra como opcional na frente para não
% quebrar os .tex que já usam a forma de dois. Visual do RG
% (relatorio-gestao.sty:246): nome em destaque, cargo ao lado. As duas
% \parbox 0,6/0,4 do book não vieram: no RG o cargo segue o nome.
\NewDocumentCommand{\CargoNosCreditos}{o m m}{%
  \par\noindent\hangindent=1em\hangafter=1\relax
  \textbf{#2}%
  \IfValueT{#1}{ \crpsp@registro{\crpsp@registro@fmt{#1}}}%
  \quad #3\par}

% ── (E) cargos alinhados: \CargoNosCreditos com coluna fixa ───────
% Terceira revisão do plan-est_2026, item 1: o book 0.5.1 tinha duas
% \parbox (0,6/0,4) e os cargos saíam alinhados na vertical. Aqui a
% mesma forma, com o tagging do `paralelos' (relatorio.sty §6): a
% pessoa é um Div; cada \parbox abre o seu parágrafo (paraon dentro,
% \par e paraoff antes de fechar). Leitura: nome → registro → cargo.
%   \crpspCreditosNome   largura da coluna do nome (com o registro)
%   \crpspCreditosVao    espaço entre as colunas
\newlength{\crpspCreditosNome}
\newlength{\crpspCreditosVao}
\AtBeginDocument{%
  \setlength{\crpspCreditosNome}{0.56\linewidth}%
  \setlength{\crpspCreditosVao}{1em}}
\newcommand{\CargoAlinhado}[3]{% #1 registro (ou vazio), #2 nome, #3 cargo
  \par
  \crpsprel@paraoff
  \crpsprel@tagbegin{Div}%
  \noindent
  \parbox[t]{\crpspCreditosNome}{%
    \raggedright\hangindent=1em\hangafter=1\relax
    \crpsprel@paraon
    \textbf{#2}%
    % registro sempre na linha de baixo, como nas duas colunas: na
    % coluna estreita ele ora cabia, ora não, e a lista ficava irregular
    \ifx\relax#1\relax\else\newline\crpsp@registro{\crpsp@registro@fmt{#1}}\fi
    \par\crpsprel@paraoff}%
  \hspace{\crpspCreditosVao}%
  \parbox[t]{\dimexpr\linewidth-\crpspCreditosNome-\crpspCreditosVao\relax}{%
    \raggedright
    \crpsprel@paraon
    #3\par\crpsprel@paraoff}%
  \par
  \crpsprel@tagend
  \crpsprel@paraon
}
\if E\variante
  \RenewDocumentCommand{\CargoNosCreditos}{o m m}{%
    \IfValueTF{#1}{\CargoAlinhado{#1}}{\CargoAlinhado{}}{#2}{#3}}
\fi

% CreditoInstitucional: o ambiente do modelo (RG e book) que envolve
% um bloco de créditos. Cada pessoa é um parágrafo, e o \parskip do
% WCAG (1,5 × a entrelinha, 21,75 pt) abria a lista demais. Aqui ele
% cai para \crpspCreditosSep. Como no sumário (comum.tex §4 do
% plan-est_2026), é exceção ao WCAG 1.4.8 numa lista de nomes, não em
% texto corrido: CONFIRMAR COM O EDITOR.
\newlength{\crpspCreditosSep}
\setlength{\crpspCreditosSep}{0.3\baselineskip}
% ⚠️ Logo depois de título, o \@afterheading da crpsp-base já
% descontou o \parskip INTEIRO, e o primeiro parágrafo daqui só devolve
% o \crpspCreditosSep: o primeiro nome encavalava no título. A diferença
% é devolvida na abertura.
\newenvironment{CreditoInstitucional}
  {\par
   \if@nobreak\vspace{\dimexpr\parskip-\crpspCreditosSep\relax}\fi
   \begingroup\setlength{\parskip}{\crpspCreditosSep}}
  {\par\endgroup}

% NomesDuasColunas: o do book (multicols, \\ vira \par), com o
% separador trocado por quebra de linha.
% O \multicolsep (acima e abaixo das colunas) punha 17,6 pt entre o
% título e o primeiro nome, contra 5,7 pt na Diretoria. Aqui ele fica
% zerado: o espaço de cima é o do hook env/multicols/before e o de
% baixo é o do título seguinte.
\newenvironment{NomesDuasColunas}{%
  \par\begingroup
  \setlength{\multicolsep}{0pt}%
  \begin{multicols}{2}%
  \begingroup
  \raggedright
  \setlength{\parskip}{\crpspCreditosSep}%
  \let\\\par
  \renewcommand{\crpsp@registro@sep}{\newline}%
}{%
  \endgroup
  \end{multicols}%
  \endgroup
}

\makeatother

\hypersetup{pdftitle={MWE — créditos com título invisível}}

\begin{document}

% ════════════════════════ página de créditos
\if B\variante\else\if C\variante\else\creditosoculto{Créditos institucionais}\fi\fi
\if B\variante\creditosaltvazio{Créditos institucionais}\fi

\creditosgrupo{XVIII Plenário (2025–2028)}

\creditossub{Diretoria}
\begin{CreditoInstitucional}
\CargoNosCreditos[06/31093]{Valéria Campinas Braunstein}{presidenta}
\CargoNosCreditos[06/149691]{Victória Soares Vidal}{vice-presidenta}
\CargoNosCreditos[06/148611]{Fabiana Macena Luiz}{secretária}
\CargoNosCreditos[06/159297]{Genildo Gomes de Sousa}{tesoureiro}
\end{CreditoInstitucional}

\creditossub{Conselheiras/os efetivas/os}
\begin{NomesDuasColunas}
  \NomeNosCreditos{Cláudia Cristina Lofrano}{06/44926}
  \NomeNosCreditos{Débora Nascimento Santos}{06/144738}
  \NomeNosCreditos{Fabiana Macena Luiz}{06/148611}
  \NomeNosCreditos{Gabriel Basílio Barbosa Costa}{06/185699}
  \NomeNosCreditos{Genildo Gomes de Sousa}{06/159297}
  \NomeNosCreditos{Luísa Thomazini de Freitas}{06/159754}
  \NomeNosCreditos{Luke Ribeiro Mazzei França Barros}{06/188231}
  \NomeNosCreditos{Marilia Capponi}{06/81224}
  \NomeNosCreditos{Natali de Souza Nascimento}{06/134212}
  \NomeNosCreditos{Paula Andréia de Carvalho Jonas}{06/62340}
\end{NomesDuasColunas}

\creditossub{Conselheiras/os suplentes (forma antiga, deprecada)}
\begin{NomesDuasColunas}
  \NomeNosCreditos{Beatris Guarita Dotta (CRP~06/143345)}
  \NomeNosCreditos{Bruna Pessenda (CRP~06/137732)}
\end{NomesDuasColunas}

\creditosgrupo{Comissões permanentes}

\creditossub{Comissão de Ética (COE)}
\begin{CreditoInstitucional}
\CargoNosCreditos{Cláudia Cristina Lofrano}{presidenta}
\end{CreditoInstitucional}

\clearpage
% ════════════════════════ segunda página: título oculto no topo,
% seguido direto de texto, para medir que nada foi empurrado
\if B\variante\else\if C\variante\else\creditosoculto{Créditos técnico-administrativos}\fi\fi
\if B\variante\creditosaltvazio{Créditos técnico-administrativos}\fi
\creditosgrupo{Gerências}
\begin{CreditoInstitucional}
\CargoNosCreditos{Vanessa Valente}{gerente de Administração e Tecnologia da Informação}
\CargoNosCreditos{Lilian Suzuki}{gerente Técnica e Política (interina)}
\CargoNosCreditos[06/110639]{Maria Carolina Vitalino Pinto Ferraz Cabau}{subcoordenadora}
\CargoNosCreditos{Adolfo Barros Benevenuto}{coordenador de Tecnologia da Informação e Comunicação}
\end{CreditoInstitucional}

\end{document}
"; done; done
%% Conferir:
%%   ~/verapdf/verapdf --format text --flavour ua2 mwe_creditos_A.pdf
%%   texlua $(kpsewhich show-pdf-tags.lua) --tree mwe_creditos_A.pdf
%%   pdftotext -bbox mwe_creditos_A.pdf - | grep -E '>(XVIII|Gerências)<'
%%
%% ⚠️ Aspas Unicode reais, nunca `` '': a linha compila com
%% \tracinglostchars=3 e a Atkinson não tem a crase.

\DocumentMetadata{
  lang        = pt-BR,
  pdfversion  = 2.0,
  pdfstandard = ua-2,
  tagging     = on,
}

\documentclass[externo]{crpsp-relatorio}

\providecommand{\variante}{A}

\usepackage{multicol}
\usepackage{pdfrender}

\makeatletter

% ── (A) título invisível com texto real ─────────────────────────
% O \@startsection da classe abre a estrutura de seção (section → H1,
% dentro de um Sect que engloba o que vem depois). O texto vai em Tr 3
% (pdfrender): não pinta, mas é conteúdo marcado — extração, cópia e
% leitor de tela o encontram.
%   - \smash\rlap: a linha do título tem altura, profundidade e
%     largura zero. Sem o \smash sobram 0,143 pt (a profundidade do
%     texto invisível).
%   - \texorpdfstring: sem ele o /T da estrutura sai
%     "TextRenderingMode=InvisibleCréditos institucionais".
%   - Depois do título, o próximo título (\creditosgrupo, \large) tem
%     de cair onde cairia no topo de uma página sem nada: linha de base
%     no \topskip. O \topskip só vale para a PRIMEIRA caixa da página,
%     e aqui a primeira é a linha invisível. A emulação:
%       \vskip-\topskip  volta a referência ao topo da mancha;
%       \prevdepth       faz a cola entre linhas valer
%                        \baselineskip(\large) − \prevdepth − h
%                        = \topskip − h, que é a cola do \topskip.
%     Medido (pdftotext -bbox): o título seguinte sai em 97,3956 pt,
%     igual à página sem título (variante C). Só com \vspace{-\parskip}
%     saía 3,7 pt acima; sem correção nenhuma, 17,9 pt abaixo.
%   ⚠️ A conta supõe que o próximo título é \large. Se o título de
%   nível 2 mudar de corpo, mudar aqui.
\newlength{\crpsp@bs@grupo}
\AtBeginDocument{{\large\global\crpsp@bs@grupo=\baselineskip}}
\newcommand{\crpsp@titulo@invisivel}[1]{%
  \textpdfrender{TextRenderingMode=Invisible}{#1}}
\newcommand{\creditosoculto}[1]{%
  \@startsection{section}{1}{0pt}{0pt}{1sp}%
    {\normalfont\fontsize{1pt}{0pt}\selectfont}%
    *{\texorpdfstring{\smash{\rlap{\crpsp@titulo@invisivel{\normalsize #1}}}}{#1}}%
  \par
  \vskip-\topskip
  \prevdepth=\dimexpr\crpsp@bs@grupo-\topskip\relax
}

% ── (B) título sem conteúdo, só /Alt ────────────────────────────
\newcommand{\creditosaltvazio}[1]{%
  \tagstructbegin{tag=H1,alt={#1}}\tagstructend}

% ── modelo do RG, com as cores e fontes da classe ────────────────
% \CreditoInstitucionalUm → nível 2 com fio; Dois → nível 3.
% (D) o fio é artefato (helper da classe, relatorio.sty §5).
\newcommand{\creditosgrupo}[1]{%
  \subsection*{#1}%
  % O \@afterheading da crpsp-base já deixou um \vspace{-\parskip} (o
  % parágrafo seguinte o devolveria). Sem devolver aqui, o fio sobe
  % 21,75 pt e sai ACIMA do título. Depois do fio o desconto é refeito,
  % para o que vier em seguida (lista, título) cair onde cairia.
  \vspace{\parskip}\vspace{2pt}%
  % em modo vertical (\hrule), não \rule num parágrafo: começar o
  % parágrafo com o tagging suspenso desbalanceia os hooks de
  % parágrafo ("begin (21) and end (24)"), um por fio.
  \crpsprel@artefatoinicio
  {\color{crpsp@cor3}\hrule height 1pt\relax}%
  \crpsprel@artefatofim
  \vspace{-\parskip}\vspace{0.4\baselineskip}%
}
% multicols logo depois de título: mesmo defeito do longtable
% (comum.tex §2 do plan-est_2026). Não abre parágrafo, não devolve o
% \parskip, e a primeira linha das colunas encavala no título.
\AddToHook{env/multicols/before}{\if@nobreak\vspace{\parskip}\fi}
\newcommand{\creditossub}[1]{\subsubsection*{#1}}

% ── nomes, cargos e registro: interface do book 0.5.1 ─────────────
% book-crpsp_acessivel.sty:781-810 (13/09/2026) tem \CargoNosCreditos
% de dois argumentos, \NomeNosCreditos de um e o NomesDuasColunas. A
% decisão de 17/09 (briefings/funcionalidade-creditos-institucionais.md)
% separa o registro do nome, para que ele não seja lido como parte do
% nome. Aqui o registro vira uma estrutura própria (Span) dentro do
% parágrafo da pessoa.
%
% O nome fica \NomeNosCreditos (decisão do editor, 24/09/2026), e a
% forma antiga continua aceita pelo mesmo nome: ver adiante.
\ExplSyntaxOn
% Fecha o MC do parágrafo, abre Span com MC próprio, e reabre o MC do
% parágrafo: o par push/pop do tagpdf existe para isto.
\cs_new_protected:Npn \crpsp_creditos_registro:n #1
  {
    \tag_mc_end_push:
    \tag_struct_begin:n { tag = Span }
    \tag_mc_begin:n { }
    #1
    \tag_mc_end:
    \tag_struct_end:
    \tag_mc_begin_pop:n { }
  }
\cs_new_eq:NN \crpsp@registro \crpsp_creditos_registro:n
\ExplSyntaxOff
% Separador entre nome e registro: espaço na largura da mancha; nas
% duas colunas, quebra de linha com recuo (aprovado pelo editor em
% 24/09: registro sempre embaixo, nunca solto pela largura da coluna).
\newcommand{\crpsp@registro@sep}{ }
\newcommand{\crpsp@registro@fmt}[1]{(CRP~#1)}

% \NomeNosCreditos{nome}{registro}: uma pessoa, um parágrafo.
%   \NomeNosCreditos{Beatris Guarita Dotta}{06/143345}
% Forma antiga (book 0.5.1), um argumento com o registro no texto —
% deprecada, com aviso:
%   \NomeNosCreditos{Beatris Guarita Dotta (CRP~06/143345)}
% As duas convivem no mesmo nome: depois do primeiro argumento, se o
% próximo token (pulando espaços) for `{', é a forma nova. O tipo `g'
% do ltcmd faria o mesmo, mas saiu do kernel ("Invalid argument type
% 'g'") e só volta com o xparse obsoleto. ⚠️ A armadilha é a mesma do
% `g': na forma antiga, se o que vem depois começar por `{', é engolido
% como registro. Nos .tex existentes o que vem depois é outro
% \NomeNosCreditos, \\, `}' ou linha em branco. Quando a forma antiga
% sair, trocar por `m m'.
\ExplSyntaxOn
\NewDocumentCommand{\NomeNosCreditos}{m}
  {
    \peek_catcode_ignore_spaces:NTF \c_group_begin_token
      { \crpsp_creditos_nome:nn {#1} }
      {
        \msg_warning:nn { crpsp-creditos } { nome-antigo }
        \crpsp_creditos_nome_antigo:n {#1}
      }
  }
\cs_new_protected:Npn \crpsp_creditos_nome:nn #1#2
  {
    \par\noindent\hangindent=1em\hangafter=1\relax
    #1 \crpsp@registro@sep \crpsp@registro { \crpsp@registro@fmt {#2} }
    \par
  }
\cs_new_protected:Npn \crpsp_creditos_nome_antigo:n #1
  { \par\noindent #1 \par }
\msg_new:nnn { crpsp-creditos } { nome-antigo }
  {
    \iow_char:N\\NomeNosCreditos~de~um~argumento~está~deprecado:~use~
    \iow_char:N\\NomeNosCreditos\{nome\}\{registro\}.
  }
\ExplSyntaxOff

% \CargoNosCreditos[registro]{nome}{cargo}: os dois obrigatórios são
% os do book 0.5.1; o registro entra como opcional na frente para não
% quebrar os .tex que já usam a forma de dois. Visual do RG
% (relatorio-gestao.sty:246): nome em destaque, cargo ao lado. As duas
% \parbox 0,6/0,4 do book não vieram: no RG o cargo segue o nome.
\NewDocumentCommand{\CargoNosCreditos}{o m m}{%
  \par\noindent\hangindent=1em\hangafter=1\relax
  \textbf{#2}%
  \IfValueT{#1}{ \crpsp@registro{\crpsp@registro@fmt{#1}}}%
  \quad #3\par}

% ── (E) cargos alinhados: \CargoNosCreditos com coluna fixa ───────
% Terceira revisão do plan-est_2026, item 1: o book 0.5.1 tinha duas
% \parbox (0,6/0,4) e os cargos saíam alinhados na vertical. Aqui a
% mesma forma, com o tagging do `paralelos' (relatorio.sty §6): a
% pessoa é um Div; cada \parbox abre o seu parágrafo (paraon dentro,
% \par e paraoff antes de fechar). Leitura: nome → registro → cargo.
%   \crpspCreditosNome   largura da coluna do nome (com o registro)
%   \crpspCreditosVao    espaço entre as colunas
\newlength{\crpspCreditosNome}
\newlength{\crpspCreditosVao}
\AtBeginDocument{%
  \setlength{\crpspCreditosNome}{0.56\linewidth}%
  \setlength{\crpspCreditosVao}{1em}}
\newcommand{\CargoAlinhado}[3]{% #1 registro (ou vazio), #2 nome, #3 cargo
  \par
  \crpsprel@paraoff
  \crpsprel@tagbegin{Div}%
  \noindent
  \parbox[t]{\crpspCreditosNome}{%
    \raggedright\hangindent=1em\hangafter=1\relax
    \crpsprel@paraon
    \textbf{#2}%
    % registro sempre na linha de baixo, como nas duas colunas: na
    % coluna estreita ele ora cabia, ora não, e a lista ficava irregular
    \ifx\relax#1\relax\else\newline\crpsp@registro{\crpsp@registro@fmt{#1}}\fi
    \par\crpsprel@paraoff}%
  \hspace{\crpspCreditosVao}%
  \parbox[t]{\dimexpr\linewidth-\crpspCreditosNome-\crpspCreditosVao\relax}{%
    \raggedright
    \crpsprel@paraon
    #3\par\crpsprel@paraoff}%
  \par
  \crpsprel@tagend
  \crpsprel@paraon
}
\if E\variante
  \RenewDocumentCommand{\CargoNosCreditos}{o m m}{%
    \IfValueTF{#1}{\CargoAlinhado{#1}}{\CargoAlinhado{}}{#2}{#3}}
\fi

% CreditoInstitucional: o ambiente do modelo (RG e book) que envolve
% um bloco de créditos. Cada pessoa é um parágrafo, e o \parskip do
% WCAG (1,5 × a entrelinha, 21,75 pt) abria a lista demais. Aqui ele
% cai para \crpspCreditosSep. Como no sumário (comum.tex §4 do
% plan-est_2026), é exceção ao WCAG 1.4.8 numa lista de nomes, não em
% texto corrido: CONFIRMAR COM O EDITOR.
\newlength{\crpspCreditosSep}
\setlength{\crpspCreditosSep}{0.3\baselineskip}
% ⚠️ Logo depois de título, o \@afterheading da crpsp-base já
% descontou o \parskip INTEIRO, e o primeiro parágrafo daqui só devolve
% o \crpspCreditosSep: o primeiro nome encavalava no título. A diferença
% é devolvida na abertura.
\newenvironment{CreditoInstitucional}
  {\par
   \if@nobreak\vspace{\dimexpr\parskip-\crpspCreditosSep\relax}\fi
   \begingroup\setlength{\parskip}{\crpspCreditosSep}}
  {\par\endgroup}

% NomesDuasColunas: o do book (multicols, \\ vira \par), com o
% separador trocado por quebra de linha.
% O \multicolsep (acima e abaixo das colunas) punha 17,6 pt entre o
% título e o primeiro nome, contra 5,7 pt na Diretoria. Aqui ele fica
% zerado: o espaço de cima é o do hook env/multicols/before e o de
% baixo é o do título seguinte.
\newenvironment{NomesDuasColunas}{%
  \par\begingroup
  \setlength{\multicolsep}{0pt}%
  \begin{multicols}{2}%
  \begingroup
  \raggedright
  \setlength{\parskip}{\crpspCreditosSep}%
  \let\\\par
  \renewcommand{\crpsp@registro@sep}{\newline}%
}{%
  \endgroup
  \end{multicols}%
  \endgroup
}

\makeatother

\hypersetup{pdftitle={MWE — créditos com título invisível}}

\begin{document}

% ════════════════════════ página de créditos
\if B\variante\else\if C\variante\else\creditosoculto{Créditos institucionais}\fi\fi
\if B\variante\creditosaltvazio{Créditos institucionais}\fi

\creditosgrupo{XVIII Plenário (2025–2028)}

\creditossub{Diretoria}
\begin{CreditoInstitucional}
\CargoNosCreditos[06/31093]{Valéria Campinas Braunstein}{presidenta}
\CargoNosCreditos[06/149691]{Victória Soares Vidal}{vice-presidenta}
\CargoNosCreditos[06/148611]{Fabiana Macena Luiz}{secretária}
\CargoNosCreditos[06/159297]{Genildo Gomes de Sousa}{tesoureiro}
\end{CreditoInstitucional}

\creditossub{Conselheiras/os efetivas/os}
\begin{NomesDuasColunas}
  \NomeNosCreditos{Cláudia Cristina Lofrano}{06/44926}
  \NomeNosCreditos{Débora Nascimento Santos}{06/144738}
  \NomeNosCreditos{Fabiana Macena Luiz}{06/148611}
  \NomeNosCreditos{Gabriel Basílio Barbosa Costa}{06/185699}
  \NomeNosCreditos{Genildo Gomes de Sousa}{06/159297}
  \NomeNosCreditos{Luísa Thomazini de Freitas}{06/159754}
  \NomeNosCreditos{Luke Ribeiro Mazzei França Barros}{06/188231}
  \NomeNosCreditos{Marilia Capponi}{06/81224}
  \NomeNosCreditos{Natali de Souza Nascimento}{06/134212}
  \NomeNosCreditos{Paula Andréia de Carvalho Jonas}{06/62340}
\end{NomesDuasColunas}

\creditossub{Conselheiras/os suplentes (forma antiga, deprecada)}
\begin{NomesDuasColunas}
  \NomeNosCreditos{Beatris Guarita Dotta (CRP~06/143345)}
  \NomeNosCreditos{Bruna Pessenda (CRP~06/137732)}
\end{NomesDuasColunas}

\creditosgrupo{Comissões permanentes}

\creditossub{Comissão de Ética (COE)}
\begin{CreditoInstitucional}
\CargoNosCreditos{Cláudia Cristina Lofrano}{presidenta}
\end{CreditoInstitucional}

\clearpage
% ════════════════════════ segunda página: título oculto no topo,
% seguido direto de texto, para medir que nada foi empurrado
\if B\variante\else\if C\variante\else\creditosoculto{Créditos técnico-administrativos}\fi\fi
\if B\variante\creditosaltvazio{Créditos técnico-administrativos}\fi
\creditosgrupo{Gerências}
\begin{CreditoInstitucional}
\CargoNosCreditos{Vanessa Valente}{gerente de Administração e Tecnologia da Informação}
\CargoNosCreditos{Lilian Suzuki}{gerente Técnica e Política (interina)}
\CargoNosCreditos[06/110639]{Maria Carolina Vitalino Pinto Ferraz Cabau}{subcoordenadora}
\CargoNosCreditos{Adolfo Barros Benevenuto}{coordenador de Tecnologia da Informação e Comunicação}
\end{CreditoInstitucional}

\end{document}
